\chapter{Background}

\section{Earth Observation metadata and STAC}

The SpatioTemporal Asset Catalog (STAC) specification defines a common representation for
geospatial assets, collections, items, properties, and links \cite{stac-spec}. Its separation of
metadata from asset URLs is useful for this thesis: the ingestion boundary validates the STAC
item, while later components consume a typed scene record. Sentinel-2 is used as the
representative optical mission; its multispectral design supports cloud assessment and
vegetation indices \cite{drusch2012sentinel}.

\section{CDSE and interoperable standards}

The Copernicus Data Space Ecosystem provides access to Copernicus data and services
\cite{cdse}. The implementation uses OAuth2 client credentials at the provider boundary and
retains the existing STAC search semantics. The design is consistent with OGC API Features
principles for interoperable feature access \cite{ogc-api-features}, while remaining small
enough for offline teaching and experiments.

\section{Provenance and reproducibility}

Reproducible research requires a record of the input, transformation context, and output.
The pipeline stores canonical JSON payloads under SHA-256 content addresses and appends
source identifiers, URIs, timestamps, and retry attempts to a manifest. Reports serialize
metrics as sorted JSON and Markdown so a benchmark can be inspected without a monitoring
service.

\section{Spatial databases and WAL}

SQLite provides a compact embedded catalog. The catalog stores scene metadata, footprints,
bounding boxes, and assets, applies indexed SQL predicates before exact geometry checks, and
uses write-ahead logging (WAL) to permit readers while a writer commits \cite{sqlite-wal}.
The implementation additionally serializes operations on a shared connection with a
process-local reentrant lock.
